\section{Essential Controls, Robustness, and Reuse}
\label{app:execution}
\subsection{Textual Guidance and Framework Realizations}
\label{app:prompt-controls}
The \emph{prompt engineering baseline}, Prompt, adds five checks to Direct: identify requirements, check conditions, verify financial arithmetic, compare eligible purchases, and review the answer. Direct already specifies joint feasibility, worst-case value, and refusal. This comparison tests the effect of additional written guidance under the same policy.

We make 200 new calls per configuration for DeepSeek low/high and Sol/Luna/Terra medium. Task, catalogue, answer schema, model setting, output cap, and deadline match the retained Direct runs. The prompts add no examples or reference answers. The native-model checklist uses the same five checks without DeepSeek's additional restriction on internal program-like representations. Both versions are retained. Comparisons across run dates do not control service-version drift.

\begin{table*}[!tbp]
\centering\small\setlength{\tabcolsep}{7pt}
\input{tables/prompt_control}
\caption{\textbf{Additional prompt engineering on the primary tasks.} Direct and Prompt are exact-decision percentages on the same 200 tasks per configuration; $\Delta$ is Prompt minus Direct in percentage points. Prompt U/R/O reports unsafe acceptance, false refusal, and other errors. Every new call remains in the denominator.}
\label{tab:prompt-control}
\end{table*}

The checklist changes exactness by $-6.5$ to $+1.0$ percentage points, with an improvement only for Terra. This result applies to the tested checklist and does not establish the effect of other prompt engineering strategies.

Set-Py implements the same framework procedurally. A generated \texttt{assess(attrs)} computes joint feasibility and worst-case utility for the fixed selector. Table~\ref{tab:four-way-controls} compares all four routes on identical tasks within each row. DeepSeek has 200 tasks. Sol, Luna, and Terra use the 40-task intersection with prior procedural runs, with one feasibility pair per template selected before this comparison. Reference decisions and attributes agree across those versions. Set compiles without the catalogue, which the other three routes receive. Both executable realizations make more exact decisions than Direct and Prompt in every row.

\begin{table*}[!tbp]
\centering\small\setlength{\tabcolsep}{6pt}
\input{tables/four_way_controls}
\caption{\textbf{Textual guidance and two implementations of the same framework.} Each cell gives exact-decision percentage above unsafe/false-refusal/other counts. All four routes use the same $n$ tasks within a configuration. The 40-task rows cover all 20 templates and use the pre-existing task intersection. There are 2,080 route outcomes on overlapping cohorts.}
\label{tab:four-way-controls}
\end{table*}

\subsection{Matched Execution Controls}
\label{app:execution-controls}
Table~\ref{tab:execution-ablation} reports safety, completion, and exactness for each realization. Procedural runs in the primary and effort studies see the catalogue. Matched-input procedural runs exclude it, as Set does. In the primary DeepSeek cohort, both realizations complete 186/200 feasible outcomes. Set records nine unsafe acceptances and Set-Py eighteen. The pooled unsafe-rate difference is $-2.25$ points (95\% template-bootstrap interval $[-6.75,+0.75]$). Relative exactness varies across the other groups despite the shared decision policy.

\begin{table*}[!tbp]
\centering\footnotesize\setlength{\tabcolsep}{4pt}
\input{tables/execution_controls}
\caption{\textbf{Risk and useful completion across framework implementations.} Arrows are Set-Py$\to$Set; $n$ is per implementation. Unsafe counts unauthorized purchases; Safe requires exact purchase and utility on feasible tasks. Effort groups cover 23 completed configurations. Matched-input controls use low/high/max, with 72 tasks per setting. Cohorts and reused Set results are not pooled across groups.}
\label{tab:execution-ablation}
\label{tab:authorization}
\end{table*}

\subsection{Requested Reasoning Settings}
\label{app:grid}
The 60-task panel combines ten historical feasibility pairs and one pair from each primary template. Figure~\ref{fig:robustness} and Table~\ref{tab:grid} cover 18 complete Sol/Terra/Luna settings. Set improves in 17, with Luna off decreasing from 25/60 to 14/60. The 23-configuration panel contributes 4,140 outcomes across three routes; Table~\ref{tab:execution-ablation} reports its grouped implementation controls.

\begin{table*}[t]
\centering\small\setlength{\tabcolsep}{5pt}
\input{tables/appendix_grid}
\caption{All complete GPT-5.6 reasoning settings: exact/$n$ and unsafe/false-refusal/other (U/R/O), with $n=60$ in every cell. The 18 configurations cover six settings each for Sol, Terra, and Luna. Settings denote requested effort, not equal or measured compute.}
\label{tab:grid}
\end{table*}

\subsection{Task Language}
\label{app:language}
The original language study evaluates the same 60 identifiers under four profiles, two languages, and three routes (1,440 outcomes). A reviewed 183-phrase translation map preserves quantities, types, events, and authorization, with system instructions held fixed. Table~\ref{tab:language} reports Direct/Set counts from these separate retained runs and the separate 200-task Astra language aggregates. The Astra updates report identical partitions, which alone cannot establish identical outputs or language invariance.

\begin{table*}[t]
\centering\small\setlength{\tabcolsep}{4pt}
\input{tables/appendix_language}
\caption{Language outcomes and failure partitions in one table. U/R/O denotes unsafe/false-refusal/other. Panels with $n=60$ and $n=200$ are identified separately. Figure~\ref{fig:robustness}b shows the DeepSeek subset.}
\label{tab:language}
\end{table*}

\subsection{Output-Budget Control}
\label{app:budget}
On 46 historically budget-sensitive Direct tasks, 92 new calls compare the original cap with 262,144 tokens under the same 1,800-second deadline. Exact decisions rise from 10/46 to 26/46, with unsafe acceptance and false refusal still observed. The original outcomes remain in their cohorts. Results on this selected subset do not estimate general failure rates.

\begin{table*}[t]
\centering\small\setlength{\tabcolsep}{5pt}
\begin{tabular}{@{}llrrrrrr@{}}
\toprule
Configuration & Method & $n$ per cap & Exact & Budget & Unsafe & Refusal & Other\\
\midrule
DS high & Direct & 21 & $4\to12$ & $11\to0$ & $0\to3$ & $6\to6$ & $0\to0$\\
DS low & Direct & 25 & $6\to14$ & $11\to1$ & $0\to1$ & $7\to7$ & $1\to2$\\
\bottomrule
\end{tabular}
\caption{Original$\to$expanded output cap for Direct; the columns partition each row's outcomes. Other denotes a remaining value error. Selection of historically budget-sensitive tasks limits prevalence estimates.}
\label{tab:budget}
\end{table*}

\subsection{Rule Reuse and Resource Accounting}
\label{app:cost}
Two fixed requirements each have eight catalogue snapshots. Direct calls the model for each snapshot, while Set compiles once per requirement. Table~\ref{tab:reuse} includes the initial compilation and all recorded input/output tokens. Request times are summed across calls and therefore do not measure parallel wall-clock latency. Reuse requires unchanged authorization, state semantics, and attribute meanings.

\begin{table*}[!tbp]
\centering\small
\begin{tabular}{llrrr}
\toprule
Model & Method & Model calls & Input + output tokens & Serial request seconds\\
\midrule
Astra max & Direct & 16 & 175,429 & 344.64\\
& Set & 2 & 22,160 & 81.69\\
\midrule
DeepSeek low & Direct & 16 & 136,345 & 351.47\\
& Set & 2 & 12,418 & 33.61\\
\bottomrule
\end{tabular}
\caption{Direct/Set reuse costs including cold compilation, over 16 catalogue evaluations per row. Both Astra routes and DeepSeek Set are exact on all snapshots; DeepSeek Direct is exact on 15/16. }
\label{tab:reuse}
\end{table*}
